\documentclass[10pt]{article}

\usepackage[utf8]{inputenc}

\usepackage[T1]{fontenc}

\usepackage{amsmath}

\usepackage{amsfonts}

\usepackage{amssymb}

\usepackage[version=4]{mhchem}

\usepackage{stmaryrd}

\usepackage{graphicx}

\usepackage[export]{adjustbox}

\graphicspath{ {./images/} }

\usepackage{bbold}



\begin{document}

ПОЛУПРАВИЛЬНЫЕ МНОГОГРАННИКИ, т е л а А p х и м е д а,- выпуклые многогранники, все грани к-рых суть правильные многоугольники, а многогранные углы конгруәнтны или симметричны. Данные о П. м. приведены в таблице, где $B$ - число вершин, $P$ - число ребер, $\Gamma$ - число граней, $\Gamma_{k}$ - число $n_{k^{-}}$

\includegraphics[max width=\textwidth, center]{2023_07_08_28610162d88fa755b7e5g-1(2)}



Полуправильны е многограннин и



\begin{center}

\includegraphics[max width=\textwidth]{2023_07_08_28610162d88fa755b7e5g-1(1)}

\end{center}



угольных гранөй, $s$ - число граней, сходящихся в каждой вершине, в том числе $s_{1} n_{1}$-угольных, $s_{2} n_{2}$-Угольных И т. д, В евклидовом пространстве $\mathbb{R}^{3}$ существует $13 \Pi$ м. [см. рис., 1-14, иногда выделяют два вида ромбокубооктаэдра (рис., 3-4), к-рые различаются тем, что верхняя часть многоугольника, состоящая из 5 квадратов и 4 правильных треугольников, повернута как целое на угол $\pi / 4]$ и две бесконечные серии - призмы (рис., 15) и антипризмы (рис., 16).



Невыпуклых (звездчатых) П. м. больше 51.



Лит.: [1] Энциклопедия әлементарной математики, кн. 4Геометрия, М.-Л., 1963; [2] Л ю с т е р ни к Л. А., Выпуклые



\includegraphics[max width=\textwidth, center]{2023_07_08_28610162d88fa755b7e5g-1}

Vielecke und Vielflache. Theorie und Geschichte, Lpz., 1900 [4] В ен нин д ж е р М., Модели многогранников, пер. с

англ., М., 1974.

А. Б. Иванов.



ПОЛУПРОСТАЯ АЛГЕБРА о т н о с и т е л ь н р а д и к а л а $r$ - алгебра, являющаяся $r$-полупростым колыцом (см. Полупростое кольцо). В нек-рых классах алгебр при подходящем выборе радикала $r$ удается описать строение П. а. (см. Классически полупростое кольцо, Альтернативные кольца и алгебры, Иорданова алгебра, Ли полупростая алгебра). ственные значения, являющиеся элементами группы $X(T)$, наз. в е с а м и п р е дс т а в л е н и я $\rho$. Ненулевые веса присоединенного представления Ad наз. к о р н я м г группы G. Оказывается, что система $\Sigma \subset X(T)$ всех корней группы $G$ является приведенной корневой системой в пространстве $E$, причем неприводимые компоненты системы $\Sigma$ - это системы корней простых замкнутых нормальных подгрупп группы $G$. Далее, $Q(\Sigma) \subseteq X(\boldsymbol{T}) \subseteq P(\Sigma)$, где $Q(\Sigma)-$ решетка радикальных весов, а $P(\Sigma)=\left\{\lambda \in E \mid \alpha^{*}(\lambda) \in \mathbb{Z}\right.$ для всех $\alpha \in \Sigma\}-$ р ші е т к а в е с о в корневой системы $\Sigma$. В случае $K=\mathbb{C}$ пространство $E$ естественно отождествляется $\quad$ вещественным подпространством $\mathrm{t}_{\mathbb{R}}^{*} \subset \mathrm{t}^{*}$, где $\mathrm{t}$ - алгебра Ли тора $T$, натянутом на дифференциалы всех характеров, а решетки в $\mathfrak{t}$, двойственные к $Q(\Sigma) \sqsubseteq X(T) \sqsubseteq P(\Sigma)$, совпадают (с точностыю до множителя 2лi) с $\Gamma_{1} \supseteq \Gamma(G) \sqsupseteq \Gamma_{0}$ (см. Ли полупростая арynna).



Основная теорема классификации утверждает, что если $G^{\prime}$ - другая П. а. г., $T^{\prime}$ - ее максимальный тор, $\Sigma^{\prime} \subset E^{\prime}$ - система корней группы $G^{\prime}$ и если существует линейное отображение $E \rightarrow E^{\prime}$, определяюпее изомор-





\end{document}